\chapter{潜能与实绩：能力包络、运行激活与闭环证据}
\label{chp:duality_potentiality_actuality}
上一章把记忆治理分解为对象、范围、操作与证据，本章转向一个更一般的问题：系统具备某种结构或算法，并不等于它在当前任务中已经表现出相应能力；一次成功输出也不等于系统拥有可迁移、可持续的能力。原有“潜能—实存”区分抓住了容量与实现不能混同这一要点，但以拓扑荷、狄拉克零模、量子内测量和边界做功直接定义意识与智能，会把数学结构、主体体验、任务表现和物理能耗压成同一个存在量。我们将其重建为能力包络、运行资格、可观测绩效和因果闭环四层模型，并分别说明意识相关功能与智能行为能够获得何种证据，避免由内部复杂度或外部动作推出本体结论。

\section{本体论定义：从结构容量到任务相关能力包络}
\label{sec:section_9823}
潜能不是系统内部某个脱离任务的标量，更不是纤维丛示性类自动赋予的“体验资格”。同一模型在数学证明、视觉控制和社会协商中的可用能力不同；同一套参数在工具断开、上下文不足或权限受限时也会表现不同。我们因此先定义任务族 $\mathcal T$、环境族 $\mathcal E$、资源集合 $\mathcal R$、接口集合 $\mathcal J$ 与安全约束集合 $\mathcal K$。系统版本 $v$ 的能力包络写成
\[
\mathcal C_v=
\{(\tau,e,r,j,k,\gamma):
\mathrm{Feasible}_v(\tau,e,r,j,k;\gamma)=1\},
\]
其中 $\tau\in\mathcal T$ 是任务，$e\in\mathcal E$ 是环境条件，$r\in\mathcal R$ 是时间、算力、带宽和外部工具预算，$j\in\mathcal J$ 是观测与行动接口，$k\in\mathcal K$ 是必须满足的边界，$\gamma$ 是预先声明的成功标准。$\mathcal C_v$ 表示经过证据支持的可行区域，而非系统在任何未来情境中“可能做到”的无限集合。

能力包络至少由五类对象共同决定。第一类是结构容量，包括参数、程序、局部结构网络、全局分布表示与外部知识；第二类是运行容量，包括当前工作窗口、并发预算、时延和故障余量；第三类是接口容量，包括传感器分辨率、工具权限、动作精度和环境反馈；第四类是学习容量，即系统在给定数据与风险预算下修订结构的范围；第五类是治理容量，包括身份、授权、审计和回滚。只用参数规模或某个拓扑指标描述潜能，会漏掉其余四类限制。

原章用陈类、A-roof 类和非平凡和乐定义意识潜力，又用目的论狄拉克算子的指标定义智能潜力。这些数学对象只有在系统确实具备相应流形、丛、联络、算子定义域和椭圆性条件时才有专门含义；即使条件成立，指标也描述解空间或拓扑不变量，不会自动成为意识或智能的量尺。我们保留“结构限制可行模式”的一般思想，把专门拓扑量降为某类模型中的特征，而不由非零拓扑不变量推出自我感、体验深度或问题求解能力。

对有限测试集，可把能力估计写成带不确定性的向量
\[
\widehat{\mathbf c}_v=
(\widehat m_1,\ldots,\widehat m_d;
\mathbf n,\mathbf b,\boldsymbol\epsilon),
\]
其中 $\widehat m_i$ 是准确率、校准、时延、资源、鲁棒性或安全性等指标，$\mathbf n$ 记录样本量与重复次数，$\mathbf b$ 记录测试分布边界，$\boldsymbol\epsilon$ 记录统计区间与测量误差。指标向量不能随意压成单一分数：高准确率可能伴随高延迟，高自主性可能伴随权限越界，低能耗也可能来自任务被简化。所谓潜能，只能在给定条件下以多维包络表达。

潜能与当前运行状态还需分离。设 $a_k$ 为时刻 $k$ 的激活与路由状态，$u_k$ 为控制和权限配置，$x_k$ 为任务相关工作态，则实际可调用子集为
\[
\mathcal C_v^{\mathrm{on}}(k)
=\Gamma(\mathcal C_v,a_k,u_k,x_k).
\]
$\Gamma$ 是筛选接口而非坍缩算子：它根据当前资源、身份和上下文选择可用能力。断开工具、冻结写权限或缩短上下文会收缩在线子集，却不必删除慢结构；恢复接口后能力可能重新出现。反过来，某个偶然成功也不能扩张整个能力包络，必须通过重复、迁移和干预测试确认。

局部结构网络与全局分布表示在这里承担不同角色。分布表示支持相似性泛化和候选生成，局部结构网络保存任务身份、对象关系、前提与约束；二者通过绑定、耦合与读出接口联合。拥有庞大分布表示而缺少来源和约束，可能生成流畅但不可靠的候选；拥有精确结构图而缺少泛化接口，则可能无法处理新表述。能力来自二者在任务条件下的协同，而不是把其中一方等同人脑某区域后推出机制相同。

验证能力包络需要三种证据。行为证据检查系统能否在留出任务上达到标准；干预证据改变资源、工具、身份或结构，观察性能是否按模型预期变化；持续性证据检查能力在时间、扰动和版本更新后是否保持。若只观察静态结构，我们得到的是实现条件；若只观察一次结果，我们得到的是事件；只有结构、干预与重复绩效共同闭合时，才能把某一区域纳入 $\mathcal C_v$。由此，“能成为什么”不再是抽象本征态库，而是一个有条件、有版本、可反驳且需要持续校准的能力承诺。
